\section{Delete the Middle Node}
\label{sec:ll-delete-middle}

\problemstatement{Given the head of a linked list, delete the middle node
(for even length, the second middle) and return the head. If the list has
one node, return null.}

\begin{patternbox}
Deleting the middle combines the slow/fast middle-finder with the need for
the node \emph{before} the middle. Advance \texttt{fast} by two but start it
one step ahead (or track \texttt{prev}) so that when \texttt{fast} reaches
the end, \texttt{slow} stops just before the middle, ready to unlink it.
\end{patternbox}

\subsection*{Stop slow one node early}

Run the tortoise--hare, but arrange the offsets so \texttt{slow} lands on the
\emph{predecessor} of the middle rather than the middle itself -- for
instance by giving \texttt{fast} a two-node head start. Then
\texttt{slow->next = slow->next->next} removes the middle node. Handle the
single-node list as a special case (return null).

\begin{ideabox}[Delete Needs the Predecessor]
Finding the middle is not enough to delete it in a singly linked list; you
need its predecessor to rewire \texttt{next}. Offsetting the fast pointer so
slow halts one node early yields that predecessor directly, avoiding a
second pass.
\end{ideabox}

\subsection*{Algorithm}

\begin{enumerate}
  \item If the list has one node, return null.
  \item \texttt{slow = head}, \texttt{fast = head->next->next}.
  \item While \texttt{fast} and \texttt{fast->next}: advance \texttt{slow}
        by one, \texttt{fast} by two.
  \item \texttt{slow->next = slow->next->next}; return the head.
\end{enumerate}

\subsection*{Dry run}

\dryrun{$1\to2\to3\to4$ (delete second middle, node $3$).}

\begin{itemize}
  \item slow$=1$, fast$=3$. Step: slow$=2$, fast$=$null.
  \item Remove $2->next$ (node $3$). Result $1\to2\to4$.
\end{itemize}

\subsection*{C++ implementation}

\begin{lstlisting}
#include <bits/stdc++.h>
using namespace std;

struct ListNode {
    int val; ListNode* next;
    ListNode(int v) : val(v), next(nullptr) {}
};

ListNode* build(const vector<int>& a) {
    ListNode dummy(0), *t = &dummy;
    for (int x : a) { t->next = new ListNode(x); t = t->next; }
    return dummy.next;
}

ListNode* deleteMiddle(ListNode* head) {
    if (!head->next) return nullptr;              // single node
    ListNode* slow = head;
    ListNode* fast = head->next->next;            // two-node head start
    while (fast && fast->next) {
        slow = slow->next;
        fast = fast->next->next;
    }
    slow->next = slow->next->next;                // unlink the middle
    return head;
}

void print(ListNode* h) { for (; h; h = h->next) cout << h->val << " "; cout << "\n"; }

int main() {
    print(deleteMiddle(build({1, 2, 3, 4, 5})));   // 1 2 4 5
    print(deleteMiddle(build({1, 2, 3, 4})));      // 1 2 4
    return 0;
}
\end{lstlisting}

\begin{complexitybox}
\textbf{Time Complexity.} $O(n)$. \textbf{Space Complexity.} $O(1)$.
\end{complexitybox}

\begin{takeawaybox}
To delete the middle, offset the fast pointer so slow halts at the middle's
predecessor -- then a single rewire removes it. Adjusting the two-pointer
offset to land on the predecessor is a common tweak to the middle-finder.
\end{takeawaybox}
